\documentclass[12pt,a4paper]{article}
\usepackage[utf8]{inputenc}
\usepackage[T2A]{fontenc}
\usepackage[russian]{babel}
\usepackage{geometry}
\usepackage{booktabs}
\usepackage{float}
\geometry{top=2cm, left=2cm, right=2cm, bottom=2cm}
\usepackage{hyperref}

\title{Теория и вопросы защиты лабораторной работы 3:\\ Линейная регрессия методами анализа данных}
\author{Ануфриев А.С. 465029}
\date{07.11.2026}

\begin{document}
\maketitle

\section{Теоретические сведения}

\subsection{Множественная линейная регрессия}
Зависимость результативного признака $y$ от $p$ факторов $x_1,\dots,x_p$ описывается линейной моделью
\begin{equation}
y = a + b_1 x_1 + \dots + b_p x_p + \varepsilon ,
\end{equation}
где $a$ --- свободный член, $b_j$ --- регрессионные коэффициенты, $\varepsilon$ --- случайная ошибка. Эмпирическая оценка:
\begin{equation}
\hat{y} = a + b_1 x_1 + \dots + b_p x_p .
\end{equation}

\subsection{Метод наименьших квадратов}
Оценки параметров получаются минимизацией суммы квадратов остатков $Q = \sum_i e_i^2 = (Y-XB)^\top(Y-XB)$. Дифференцирование по $B$ даёт нормальное уравнение
\begin{equation}
B = (X^\top X)^{-1} X^\top Y ,
\end{equation}
где матрица $X$ содержит первый столбец единиц (для свободного члена) и столбцы значений факторов, а $Y$ --- вектор наблюдений целевой переменной.

\subsection{Предположения МНК (теорема Гаусса---Маркова)}
\begin{enumerate}
\item модель специфицирована верно (линейная зависимость);
\item между факторами нет строгой линейной зависимости (отсутствие мультиколлинеарности);
\item ошибки случайны, математическое ожидание равно нулю, дисперсия постоянна (гомоскедастичность);
\item ошибки не коррелированы между собой.
\end{enumerate}
При выполнении этих условий оценки МНК являются наилучшими линейными несмещёнными оценками.

\subsection{Коэффициент детерминации}
\begin{equation}
R^2 = 1 - \frac{\sum_i (y_i - \hat{y}_i)^2}{\sum_i (y_i - \bar{y})^2} ,
\end{equation}
показывает долю дисперсии целевой переменной, объясняемую моделью. Чем ближе $R^2$ к единице, тем лучше качество модели.

\section{Вероятные вопросы защиты с ответами}

\begin{enumerate}
\item \textbf{Что такое метод наименьших квадратов и в чём его суть?}\\
МНК --- способ оценки параметров регрессионной модели, минимизирующий сумму квадратов отклонений наблюдаемых значений зависимой переменной от значений, предсказанных моделью. Аналитически решение задаётся формулой $B = (X^\top X)^{-1}X^\top Y$.

\item \textbf{Зачем нужен первый столбец единиц в матрице $X$?}\\
Столбец единиц соответствует свободному члену $a$ модели --- значению $y$, когда все факторы равны нулю. Он позволяет смещать регрессионную плоскость по оси $y$.

\item \textbf{Как вы проверяли отсутствие мультиколлинеарности?}\\
В исходных данных категориальных признаков нет, все признаки числовые. В третьей модели объединены $total\_rooms$, $households$ и их отношение $rooms\_per\_household = total\_rooms/household$. Это создаёт строгую функциональную связь, что проявляется как мультиколлинеарность: модель 3 не превосходит модель 2, а отдельные коэффициенты нестабильны.

\item \textbf{Что означает коэффициент детерминации $R^2$?}\\
Доля дисперсии целевой переменной, объясняемая регрессионной моделью. $R^2 = 1$ --- идеальное описание, $R^2 = 0$ --- модель не лучше среднего значения. В работе лучше модель 2 с $R^2 \approx 0.54$.

\item \textbf{Зачем нужно разделение на обучающую и тестовую выборки?}\\
Чтобы оценить качество модели на данных, которые не использовались при обучении, и выявить переобучение. Деление 80\%/20\% (13 600/3 400 записей).

\item \textbf{Что делает нормировка (стандартизация) признаков?}\\
Приводит признаки к сравнимому масштабу (вычет среднего, деление на стандартное отклонение), что упрощает интерпретацию коэффициентов и устойчивость вычислений. Целевая переменная не нормировалась.

\item \textbf{Какой признак оказал наибольшее влияние на стоимость жилья?}\\
Доход населения ($median\_income$) --- сильнейший отдельный фактор ($R^2 \approx 0.50$ в модели 1). Добавление возраста дома повысило качество до $R^2 \approx 0.54$ (модель 2).

\item \textbf{Почему синтетический признак не улучшил модель 3?}\\
Признак $rooms\_per\_household$ линейно зависит от $total\_rooms$ и $households$, поэтому при их совместном включении возникает мультиколлинеарность, и качество не растёт ($R^2 \approx 0.52 < 0.54$ модели 2).

\item \textbf{Какие статистики были вычислены по датасету?}\\
Количество записей (17 000), среднее, стандартное отклонение, минимум, максимум, квантили 25\%, 50\% (медиана), 75\% для всех девяти признаков.

\item \textbf{Какие библиотеки использовались в реализации?}\\
Только NumPy и Pandas (загрузка данных, работа с массивами); коэффициенты регрессии вычислены собственным кодом через матричные операции NumPy, готовые средства линейной регрессии не применялись.
\end{enumerate}

\end{document}
